\section{Linearization, regular representation, and group algebra}

Let \(k\) be a field and let \((\Omega,\rho)\) be a set-representation of \(G\). Thus
\[
    \rho:G\longrightarrow \operatorname{Sym}(\Omega)
\]
is a group homomorphism.

Denote by
\[
    k\Omega
\]
the \(k\)-vector space with basis \(\Omega\). That is, every element of \(k\Omega\) is a
finite \(k\)-linear combination
\[
    \sum_{\omega\in \Omega} a_\omega \omega,
    \qquad a_\omega\in k.
\]

The action of \(G\) on \(\Omega\) induces a \(k\)-linear representation
\[
    k\rho:G\longrightarrow \operatorname{GL}(k\Omega)
\]
defined on basis elements by
\[
    (k\rho)(g)(\omega)=\rho(g)(\omega)=g\cdot \omega,
    \qquad g\in G,\ \omega\in \Omega,
\]
and extended \(k\)-linearly.

\begin{lemma}
    The map
    \[
        k\rho:G\longrightarrow \operatorname{GL}(k\Omega)
    \]
    is a group homomorphism.
\end{lemma}

\begin{proof}
    For each \(g\in G\), the map \(\rho(g):\Omega\to\Omega\) is a bijection. Hence
    its \(k\)-linear extension
    \[
        (k\rho)(g):k\Omega\to k\Omega
    \]
    is an invertible \(k\)-linear map.

    For \(g,h\in G\) and \(\omega\in \Omega\), we have
    \[
        (k\rho)(gh)(\omega)
        =
        (gh)\cdot \omega
        =
        g\cdot(h\cdot \omega)
        =
        (k\rho)(g)\bigl((k\rho)(h)(\omega)\bigr).
    \]
    Since the basis elements \(\omega\in\Omega\) span \(k\Omega\), it follows that
    \[
        (k\rho)(gh)=(k\rho)(g)\circ (k\rho)(h).
    \]
    Also,
    \[
        (k\rho)(1_G)=\operatorname{id}_{k\Omega}.
    \]
    Therefore \(k\rho\) is a group homomorphism.
\end{proof}

\begin{definition}[Linearization of a \(G\)-set]
    The \(k\)-linear representation
    \[
        (k\Omega,k\rho)
    \]
    is called the linearization of the \(G\)-set \((\Omega,\rho)\).
\end{definition}

Let
\[
    f:(\Omega,\rho)\longrightarrow (\Lambda,\sigma)
\]
be a morphism of \(G\)-sets. Thus \(f:\Omega\to \Lambda\) is a map satisfying
\[
    f(g\cdot \omega)=g\cdot f(\omega)
\]
for all \(g\in G\) and \(\omega\in\Omega\).

Define a \(k\)-linear map
\[
    kf:k\Omega\longrightarrow k\Lambda
\]
by
\[
    kf(\omega)=f(\omega),
    \qquad \omega\in\Omega,
\]
and extend \(k\)-linearly.

Then \(kf\) is a morphism of \(k\)-linear representations. Indeed, for \(g\in G\) and
\(\omega\in\Omega\),
\[
    kf\bigl((k\rho)(g)(\omega)\bigr)
    =
    kf(g\cdot\omega)
    =
    f(g\cdot\omega)
    =
    g\cdot f(\omega)
    =
    (k\sigma)(g)(kf(\omega)).
\]
Hence
\[
    kf\circ (k\rho)(g)=(k\sigma)(g)\circ kf.
\]

Therefore the construction
\[
    (\Omega,\rho)\longmapsto (k\Omega,k\rho)
\]
induces a functor
\[
    k[-]:
    \operatorname{Rep}(G,\operatorname{Set})
    \longrightarrow
    \operatorname{Rep}(G,\operatorname{vect}_k).
\]


Let \((\Omega,\rho)\) and \((\Lambda,\sigma)\) be two \(G\)-sets.

\begin{proposition}
    There is a natural isomorphism of \(k\)-linear representations
    \[
        k(\Omega\sqcup \Lambda)
        \cong
        k\Omega\oplus k\Lambda.
    \]
\end{proposition}

\begin{proof}
    The disjoint union \(\Omega\sqcup\Lambda\) has basis elements coming either from
    \(\Omega\) or from \(\Lambda\). Hence the assignment
    \[
        \omega\longmapsto (\omega,0),
        \qquad
        \lambda\longmapsto (0,\lambda)
    \]
    induces a \(k\)-linear isomorphism
    \[
        k(\Omega\sqcup\Lambda)
        \cong
        k\Omega\oplus k\Lambda.
    \]
    This isomorphism is \(G\)-equivariant because \(G\) acts on each summand by the
    corresponding action.
\end{proof}

\begin{proposition}
    There is a natural isomorphism of \(k\)-linear representations
    \[
        k(\Omega\times \Lambda)
        \cong
        k\Omega\otimes_k k\Lambda.
    \]
\end{proposition}

\begin{proof}
    Define
    \[
        \Phi:k(\Omega\times\Lambda)\longrightarrow k\Omega\otimes_k k\Lambda
    \]
    on basis elements by
    \[
        \Phi(\omega,\lambda)=\omega\otimes \lambda.
    \]
    Since the elements \((\omega,\lambda)\) form a basis of \(k(\Omega\times\Lambda)\),
    this defines a \(k\)-linear map.

    The elements
    \[
        \omega\otimes\lambda,
        \qquad \omega\in\Omega,\ \lambda\in\Lambda,
    \]
    form a basis of \(k\Omega\otimes_k k\Lambda\), so \(\Phi\) is a \(k\)-linear
    isomorphism.

    It is \(G\)-equivariant because
    \[
        \Phi\bigl(g\cdot(\omega,\lambda)\bigr)
        =
        \Phi(g\cdot\omega,g\cdot\lambda)
        =
        (g\cdot\omega)\otimes(g\cdot\lambda),
    \]
    while
    \[
        g\cdot\Phi(\omega,\lambda)
        =
        g\cdot(\omega\otimes\lambda)
        =
        (g\cdot\omega)\otimes(g\cdot\lambda).
    \]
    Hence
    \[
        k(\Omega\times \Lambda)
        \cong
        k\Omega\otimes_k k\Lambda
    \]
    as \(k\)-linear representations of \(G\).
\end{proof}

\begin{remark}
    Thus linearization transforms disjoint unions of \(G\)-sets into direct sums of
    \(k\)-linear representations, and transforms products of \(G\)-sets into tensor
    products of \(k\)-linear representations.
\end{remark}

\begin{remark}
    Even if \((\Omega,\rho)\) is transitive as a \(G\)-set and
    \[
        |\Omega|>1,
    \]
    the linearized representation
    \[
        (k\Omega,k\rho)
    \]
    need not be irreducible.

    Indeed, define
    \[
        S_\Omega:=\sum_{\omega\in\Omega}\omega\in k\Omega.
    \]
    Since \(\Omega\) is a basis of \(k\Omega\), we have
    \[
        S_\Omega\neq 0.
    \]
    For every \(g\in G\),
    \[
        g\cdot S_\Omega
        =
        \sum_{\omega\in\Omega} g\cdot \omega.
    \]
    Since \(g:\Omega\to\Omega\) is a bijection, the set
    \[
        \{g\cdot \omega\mid \omega\in\Omega\}
    \]
    is again \(\Omega\). Hence
    \[
        g\cdot S_\Omega=S_\Omega.
    \]
    Therefore the line
    \[
        kS_\Omega
    \]
    is \(G\)-stable. If \(|\Omega|>1\), this is a nonzero proper subrepresentation
    of \(k\Omega\). Hence \(k\Omega\) is not irreducible.
\end{remark}


\begin{definition}[Left regular representation]
    Let \(G\) act on itself by left multiplication:
    \[
        g\cdot x:=gx,
        \qquad g,x\in G.
    \]
    The linearization of this \(G\)-set is the representation
    \[
        kG
    \]
    with basis \(G\), where \(G\) acts by
    \[
        g\cdot x=gx.
    \]
    This representation is called the left regular representation of \(G\).
\end{definition}

Let
\[
    S_G:=\sum_{g\in G}g\in kG.
\]
Then the line
\[
    kS_G
\]
is \(G\)-stable, since
\[
    h\cdot S_G
    =
    \sum_{g\in G} hg
    =
    \sum_{g\in G} g
    =
    S_G
\]
for every \(h\in G\).

\begin{lemma}
    Assume that \(k\) has characteristic \(p>0\) and that
    \[
        p\mid |G|.
    \]
    Then the \(G\)-stable line
    \[
        kS_G\subseteq kG
    \]
    has no \(G\)-stable complement in the left regular representation \(kG\).
\end{lemma}

\begin{proof}
    Suppose, for contradiction, that there exists a \(G\)-stable subspace \(V'\subseteq kG\)
    such that
    \[
        kG=kS_G\oplus V'.
    \]
    Let
    \[
        \pi:kG\longrightarrow kS_G
    \]
    be the projection with kernel \(V'\). Thus
    \[
        \pi|_{kS_G}=\operatorname{id}_{kS_G}.
    \]
    Since both \(kS_G\) and \(V'\) are \(G\)-stable, the projection \(\pi\) is
    \(G\)-equivariant:
    \[
        \pi(gx)=g\pi(x),
        \qquad g\in G,\ x\in kG.
    \]

    Write
    \[
        \pi(1)=\lambda S_G
    \]
    for some \(\lambda\in k\). Then, for every \(g\in G\),
    \[
        \pi(g)=\pi(g\cdot 1)=g\cdot \pi(1)=g\cdot(\lambda S_G)=\lambda S_G.
    \]
    Therefore
    \[
        \pi(S_G)
        =
        \pi\left(\sum_{g\in G}g\right)
        =
        \sum_{g\in G}\pi(g)
        =
        \sum_{g\in G}\lambda S_G
        =
        |G|\lambda S_G.
    \]
    Since \(\operatorname{char}(k)=p\) and \(p\mid |G|\), we have
    \[
        |G|=0
        \qquad\text{in } k.
    \]
    Hence
    \[
        \pi(S_G)=0.
    \]
    But \(S_G\in kS_G\) and \(\pi\) restricts to the identity on \(kS_G\), so
    \[
        \pi(S_G)=S_G.
    \]
    This is a contradiction, because \(S_G\neq 0\). Therefore \(kS_G\) has no
    \(G\)-stable complement in \(kG\).
\end{proof}

Let \(k\) be a field, and let
\[
    (M,\rho),\qquad (N,\sigma)
\]
be two \(k\)-linear representations of a group \(G\).

The group \(G\) acts on the \(k\)-vector space
\[
    \operatorname{Hom}_k(M,N)
\]
by
\[
    g\cdot \alpha
    :=
    \sigma(g)\circ \alpha\circ \rho(g)^{-1},
    \qquad
    g\in G,\quad \alpha\in \operatorname{Hom}_k(M,N).
\]
Equivalently, the following diagram commutes:
\[
\begin{tikzcd}
    M \arrow[r,"\alpha"] \arrow[d,"\rho(g)"']
        & N \arrow[d,"\sigma(g)"] \\
    M \arrow[r,"g\cdot \alpha"']
        & N .
\end{tikzcd}
\]

\begin{lemma}
    The formula
    \[
        g\cdot \alpha
        =
        \sigma(g)\circ \alpha\circ \rho(g)^{-1}
    \]
    defines a \(k\)-linear representation of \(G\) on
    \(\operatorname{Hom}_k(M,N)\).
\end{lemma}

\begin{proof}
    For each \(g\in G\), the map
    \[
        \alpha\longmapsto \sigma(g)\circ \alpha\circ \rho(g)^{-1}
    \]
    is \(k\)-linear. Moreover, for \(g,h\in G\),
    \[
    \begin{aligned}
        g\cdot (h\cdot \alpha)
        &=
        \sigma(g)\circ
        \bigl(\sigma(h)\circ \alpha\circ \rho(h)^{-1}\bigr)
        \circ \rho(g)^{-1}  \\
        &=
        \sigma(gh)\circ \alpha\circ
        \bigl(\rho(h)^{-1}\circ \rho(g)^{-1}\bigr) \\
        &=
        \sigma(gh)\circ \alpha\circ \rho(gh)^{-1} \\
        &=
        (gh)\cdot \alpha .
    \end{aligned}
    \]
    Also \(1_G\cdot \alpha=\alpha\). Hence this is a representation of \(G\).
\end{proof}

Taking \(N=k\) with the trivial representation, we obtain a natural representation on
the dual vector space
\[
    M^*:=\operatorname{Hom}_k(M,k).
\]

\begin{definition}[Contragredient representation]
    Let \((M,\rho)\) be a \(k\)-linear representation of \(G\). Its contragredient
    representation is the representation on
    \[
        M^*=\operatorname{Hom}_k(M,k)
    \]
    defined by
    \[
        (g\cdot \alpha)(m)
        =
        \alpha\bigl(\rho(g)^{-1}m\bigr),
        \qquad
        g\in G,\ \alpha\in M^*,\ m\in M.
    \]
    Equivalently,
    \[
        g\cdot \alpha=\alpha\circ \rho(g)^{-1}.
    \]
\end{definition}

The tensor product
\[
    M\otimes_k N
\]
also carries a natural \(G\)-representation.

\begin{definition}[Tensor product representation]
    The tensor product of \((M,\rho)\) and \((N,\sigma)\) is
    \[
        (M,\rho)\otimes (N,\sigma)
        :=
        (M\otimes_k N,\rho\otimes\sigma),
    \]
    where
    \[
        (\rho\otimes\sigma)(g)(m\otimes n)
        :=
        \rho(g)m\otimes \sigma(g)n.
    \]
    In action notation,
    \[
        g\cdot(m\otimes n)=(g\cdot m)\otimes(g\cdot n).
    \]
\end{definition}

A representation of degree \(1\) is the same as a group homomorphism
\[
    \rho:G\longrightarrow k^\times.
\]
Indeed, if \(M=k\), then
\[
    \operatorname{GL}(M)=\operatorname{GL}_1(k)\cong k^\times.
\]

\begin{definition}
    A one-dimensional representation of \(G\) over \(k\) is a pair
    \[
        (k,\rho),
    \]
    where
    \[
        \rho:G\longrightarrow k^\times
    \]
    is a group homomorphism.
\end{definition}

\begin{remark}
    Since \(k^\times\) is abelian, every homomorphism
    \[
        \rho:G\longrightarrow k^\times
    \]
    kills the commutator subgroup \([G,G]\). Hence \(\rho\) factors uniquely
    through the abelianization:
    \[
        G/[G,G]\longrightarrow k^\times.
    \]
    Therefore one-dimensional representations of \(G\) over \(k\) are naturally
    the same as group homomorphisms
    \[
        G/[G,G]\longrightarrow k^\times.
    \]
\end{remark}

\begin{example}[One-dimensional representations of \(D_8\)]
    Let
    \[
        G=D_8=\langle \sigma,\tau\mid \sigma^2=\tau^2=(\sigma\tau)^4=1\rangle.
    \]
    Then
    \[
        [G,G]=Z(G)=\langle(\sigma\tau)^2\rangle,
    \]
    and hence
    \[
        G/[G,G]\cong \mathbb{Z}_2\times \mathbb{Z}_2.
    \]
    Therefore, over a field \(k\) with \(\operatorname{char}(k)\neq 2\), the group
    \(G\) has four one-dimensional representations:
    \[
    \begin{array}{c|cc}
        & \sigma & \tau \\ \hline
        \rho_1 & 1 & 1 \\
        \rho_2 & 1 & -1 \\
        \rho_3 & -1 & 1 \\
        \rho_4 & -1 & -1
    \end{array}
    \]
\end{example}

Let \((M,\rho)\) be a \(k\)-linear representation of \(G\), and let
\[
    (k,\sigma)
\]
be a one-dimensional representation of \(G\). Then
\[
    M\otimes_k k\cong M
\]
as \(k\)-vector spaces. Under this identification, the tensor product representation
\[
    (M,\rho)\otimes(k,\sigma)
\]
corresponds to the representation
\[
    (M,\rho\otimes\sigma),
\]
where
\[
    (\rho\otimes\sigma)(g)=\sigma(g)\rho(g).
\]

\begin{lemma}
    Let \((S,\rho)\) be an irreducible \(k\)-linear representation of \(G\), and
    let \((k,\sigma)\) be a one-dimensional representation of \(G\). Then
    \[
        (S,\rho)\otimes(k,\sigma)
    \]
    is irreducible.
\end{lemma}

\begin{proof}
    Under the identification \(S\otimes_k k\cong S\), the tensor product
    representation is given by
    \[
        g\longmapsto \sigma(g)\rho(g).
    \]
    Let \(W\subseteq S\) be a nonzero subspace stable under this twisted action.
    For \(w\in W\) and \(g\in G\), we have
    \[
        \sigma(g)\rho(g)w\in W.
    \]
    Since \(\sigma(g)\in k^\times\), multiplication by \(\sigma(g)\) is invertible.
    Hence
    \[
        \rho(g)w\in W.
    \]
    Thus \(W\) is stable under the original action \(\rho\). Since \((S,\rho)\) is
    irreducible, we must have
    \[
        W=S.
    \]
    Therefore the twisted representation is irreducible.
\end{proof}

Let \(G\) be a finite group and \(k\) a field.

\begin{definition}[Group algebra]
    The group algebra \(kG\) is the \(k\)-vector space with basis \(G\):
    \[
        kG
        =
        \left\{
            \sum_{g\in G} x_g g
            \ \middle|\ 
            x_g\in k
        \right\}.
    \]
    Multiplication is defined by extending the group multiplication \(k\)-bilinearly.
    Thus, if
    \[
        x=\sum_{s\in G}x_s s,
        \qquad
        y=\sum_{t\in G}y_t t,
    \]
    then
    \[
        xy
        =
        \sum_{s,t\in G}x_s y_t\, st.
    \]
    Equivalently,
    \[
        xy
        =
        \sum_{g\in G}
        \left(
            \sum_{s\in G}x_s y_{s^{-1}g}
        \right)g.
    \]
\end{definition}

\begin{proposition}
    The group algebra \(kG\) is a \(k\)-algebra.
\end{proposition}

\begin{proof}
    The multiplication on \(kG\) is defined by \(k\)-bilinear extension of the
    multiplication in \(G\). Since multiplication in \(G\) is associative, the
    induced multiplication on \(kG\) is associative. The identity element is
    \(1_G\in G\subseteq kG\). Hence \(kG\) is a unital associative \(k\)-algebra.
\end{proof}

\begin{theorem}
    The category of \(k\)-linear representations of \(G\) is naturally equivalent
    to the category of left \(kG\)-modules:
    \[
        \operatorname{Rep}(G,\operatorname{vect}_k)
        \simeq
        kG\text{-}\operatorname{mod}.
    \]
\end{theorem}

\begin{proof}
    Let \((M,\rho)\) be a \(k\)-linear representation of \(G\). Define an action of
    \(kG\) on \(M\) by
    \[
        \left(\sum_{g\in G} a_g g\right)m
        :=
        \sum_{g\in G} a_g\,\rho(g)(m).
    \]
    This makes \(M\) into a left \(kG\)-module.

    Conversely, if \(M\) is a left \(kG\)-module, then each \(g\in G\subseteq kG\)
    acts on \(M\) by a \(k\)-linear automorphism, since \(g^{-1}\in G\) acts as its
    inverse. Hence we obtain a group homomorphism
    \[
        \rho:G\longrightarrow \operatorname{GL}(M),
        \qquad
        \rho(g)(m)=g m.
    \]
    This gives a \(k\)-linear representation of \(G\).

    These two constructions are inverse to each other, and morphisms are exactly
    the \(k\)-linear maps commuting with the \(G\)-action, equivalently the
    \(kG\)-module homomorphisms. Therefore the two categories are naturally
    equivalent.
\end{proof}


